\documentclass[11pt]{article}

\usepackage{amsmath,amssymb,amstext,amsthm, standalone}
\usepackage{geometry}
\usepackage{fancyhdr}
% \usepackage[pdftex]{graphicx}
\usepackage{xcolor}
\usepackage{tikz}
\usepackage{amssymb}

\geometry{
  verbose,
  dvips,
  width=430pt, marginparsep=0pt, marginparwidth=0pt,
  top=90pt, headheight=12pt, headsep=20pt, footskip=30pt, bottom=80pt}

\setlength{\parskip}{\medskipamount}
\setlength{\parindent}{0pt}

\linespread{1.05}

\newtheorem{theorem}{Theorem}
\newtheorem{lemma}[theorem]{Lemma}
\newtheorem{proposition}[theorem]{Proposition}
\newtheorem{corollary}[theorem]{Corollary}
\newtheorem{conjecture}[theorem]{Conjecture}
\newtheorem{obs}{Observation}
\newtheorem{clm}{Claim}
\newtheorem{ques}{Question}
\newtheorem{problem}{Problem}

\newtheorem{ex}{Example}


\newcommand{\spc}{\hspace{0.08in}}
\newcommand{\vs}{\vspace*{.1in}}
\newcommand{\E}{\mathbb{E}}
\newcommand{\Prob}{\mathbb{P}}
\newcommand{\edit}[1]{\emph{\textcolor{blue}{#1}}}


\renewenvironment{proof}{{\noindent \textbf{\textit{Proof.}}}}
{\hfill $\Box$ \vspace*{0.1in}}
\newenvironment{proofc}{{\noindent \textit{Proof of Claim.}}}
{\hfill $\Box$}


\begin{document}
\begin{huge}
    \begin{center}
        Checkpoint
    \end{center}
\end{huge}
\begin{large}\begin{center}\bf 14.6: Directional Derivatives \end{center}
\end{large}

\vs\vs

\begin{problem}
If we want to find the directional derivative in the direction of $\vec{v} = \langle 19, -2 \rangle$ what is wrong with $\vec{v}$?
\end{problem}

\begin{problem}
What is one way to interpret the formula for a directional derivative in the direction of $\hat{u} = \langle u_1, u_2 \rangle$ which is $\mathcal{D}_{\hat{u}}f(x,y) = f_xu_1+f_yu_2$?
\end{problem}

\begin{problem}
If we package our partial in a vector known as the gradient $\langle f_x, f_x \rangle = \nabla f(x,y)$ what does $\nabla f(x,y)$ and $|\nabla f(x,y)|$ tell us?
\end{problem}


\begin{problem}
If $f_x = x+e^y$ and $f_y = \ln{(x^2y)}$ find the directional derivative in the direction that is made by an angle of $\pi/3$ from the $x$-axis.
\end{problem}





\newpage

\section{Answers}

\begin{enumerate}
    \item It is not a unit vector since $|\vec{v}| = \sqrt{365}$
    \item A weighted average of our partials. We can think of $u_1$ and $u_2$ being the percentage we weight our partial in the $x$ and partial in the $y$ direction.
    \item the direction of steepest ascent from any point and the slope of the steepest ascent.
    \item $(x+e^y)(\frac{1}{2})+(\ln{(x^2y)})(\frac{\sqrt{3}}{2})$
\end{enumerate}



\end{document} 